% Options for packages loaded elsewhere
\PassOptionsToPackage{unicode}{hyperref}
\PassOptionsToPackage{hyphens}{url}
\PassOptionsToPackage{dvipsnames,svgnames,x11names}{xcolor}
%
\documentclass[
  10pt,
  a4paper,
  DIV=11]{scrartcl}

\usepackage{amsmath,amssymb}
\usepackage{iftex}
\ifPDFTeX
  \usepackage[T1]{fontenc}
  \usepackage[utf8]{inputenc}
  \usepackage{textcomp} % provide euro and other symbols
\else % if luatex or xetex
  \usepackage{unicode-math}
  \defaultfontfeatures{Scale=MatchLowercase}
  \defaultfontfeatures[\rmfamily]{Ligatures=TeX,Scale=1}
\fi
\usepackage{lmodern}
\ifPDFTeX\else  
    % xetex/luatex font selection
    \setmainfont[]{Arial}
\fi
% Use upquote if available, for straight quotes in verbatim environments
\IfFileExists{upquote.sty}{\usepackage{upquote}}{}
\IfFileExists{microtype.sty}{% use microtype if available
  \usepackage[]{microtype}
  \UseMicrotypeSet[protrusion]{basicmath} % disable protrusion for tt fonts
}{}
\makeatletter
\@ifundefined{KOMAClassName}{% if non-KOMA class
  \IfFileExists{parskip.sty}{%
    \usepackage{parskip}
  }{% else
    \setlength{\parindent}{0pt}
    \setlength{\parskip}{6pt plus 2pt minus 1pt}}
}{% if KOMA class
  \KOMAoptions{parskip=half}}
\makeatother
\usepackage{xcolor}
\usepackage[top=30mm,left=20mm,heightrounded]{geometry}
\setlength{\emergencystretch}{3em} % prevent overfull lines
\setcounter{secnumdepth}{3}
% Make \paragraph and \subparagraph free-standing
\makeatletter
\ifx\paragraph\undefined\else
  \let\oldparagraph\paragraph
  \renewcommand{\paragraph}{
    \@ifstar
      \xxxParagraphStar
      \xxxParagraphNoStar
  }
  \newcommand{\xxxParagraphStar}[1]{\oldparagraph*{#1}\mbox{}}
  \newcommand{\xxxParagraphNoStar}[1]{\oldparagraph{#1}\mbox{}}
\fi
\ifx\subparagraph\undefined\else
  \let\oldsubparagraph\subparagraph
  \renewcommand{\subparagraph}{
    \@ifstar
      \xxxSubParagraphStar
      \xxxSubParagraphNoStar
  }
  \newcommand{\xxxSubParagraphStar}[1]{\oldsubparagraph*{#1}\mbox{}}
  \newcommand{\xxxSubParagraphNoStar}[1]{\oldsubparagraph{#1}\mbox{}}
\fi
\makeatother


\providecommand{\tightlist}{%
  \setlength{\itemsep}{0pt}\setlength{\parskip}{0pt}}\usepackage{longtable,booktabs,array}
\usepackage{calc} % for calculating minipage widths
% Correct order of tables after \paragraph or \subparagraph
\usepackage{etoolbox}
\makeatletter
\patchcmd\longtable{\par}{\if@noskipsec\mbox{}\fi\par}{}{}
\makeatother
% Allow footnotes in longtable head/foot
\IfFileExists{footnotehyper.sty}{\usepackage{footnotehyper}}{\usepackage{footnote}}
\makesavenoteenv{longtable}
\usepackage{graphicx}
\makeatletter
\newsavebox\pandoc@box
\newcommand*\pandocbounded[1]{% scales image to fit in text height/width
  \sbox\pandoc@box{#1}%
  \Gscale@div\@tempa{\textheight}{\dimexpr\ht\pandoc@box+\dp\pandoc@box\relax}%
  \Gscale@div\@tempb{\linewidth}{\wd\pandoc@box}%
  \ifdim\@tempb\p@<\@tempa\p@\let\@tempa\@tempb\fi% select the smaller of both
  \ifdim\@tempa\p@<\p@\scalebox{\@tempa}{\usebox\pandoc@box}%
  \else\usebox{\pandoc@box}%
  \fi%
}
% Set default figure placement to htbp
\def\fps@figure{htbp}
\makeatother
% definitions for citeproc citations
\NewDocumentCommand\citeproctext{}{}
\NewDocumentCommand\citeproc{mm}{%
  \begingroup\def\citeproctext{#2}\cite{#1}\endgroup}
\makeatletter
 % allow citations to break across lines
 \let\@cite@ofmt\@firstofone
 % avoid brackets around text for \cite:
 \def\@biblabel#1{}
 \def\@cite#1#2{{#1\if@tempswa , #2\fi}}
\makeatother
\newlength{\cslhangindent}
\setlength{\cslhangindent}{1.5em}
\newlength{\csllabelwidth}
\setlength{\csllabelwidth}{3em}
\newenvironment{CSLReferences}[2] % #1 hanging-indent, #2 entry-spacing
 {\begin{list}{}{%
  \setlength{\itemindent}{0pt}
  \setlength{\leftmargin}{0pt}
  \setlength{\parsep}{0pt}
  % turn on hanging indent if param 1 is 1
  \ifodd #1
   \setlength{\leftmargin}{\cslhangindent}
   \setlength{\itemindent}{-1\cslhangindent}
  \fi
  % set entry spacing
  \setlength{\itemsep}{#2\baselineskip}}}
 {\end{list}}
\usepackage{calc}
\newcommand{\CSLBlock}[1]{\hfill\break\parbox[t]{\linewidth}{\strut\ignorespaces#1\strut}}
\newcommand{\CSLLeftMargin}[1]{\parbox[t]{\csllabelwidth}{\strut#1\strut}}
\newcommand{\CSLRightInline}[1]{\parbox[t]{\linewidth - \csllabelwidth}{\strut#1\strut}}
\newcommand{\CSLIndent}[1]{\hspace{\cslhangindent}#1}

%% load packages
\usepackage{sectsty}
\usepackage{fontspec}
\usepackage{scrlayer-scrpage}
\usepackage{fontspec}

% Define title formatting

%% Set font sizes for sections and subsections
\sectionfont{\fontsize{11}{16}\selectfont\bfseries} % Arial 14pt bold
\subsectionfont{\fontsize{10}{14}\selectfont} % Arial 12pt

\lehead[test1]{\pagemark}
\cehead[]{}
\rehead[test 2]{\leftmark}
\lohead[{\fontsize{8}{10}\selectfont\upshape\textbf{Institut Sekundarstufe I}} \\
\fontsize{8}{10}\selectfont\upshape Fabrikstrasse 8, CH-3012 Bern \\
\fontsize{8}{10}\selectfont\upshape T +41 31 309 21 15, contactdesk@phbern.ch, www.phbern.ch
]{\rightmark}
\rohead[PHBern]{}
\KOMAoption{captions}{tableheading}
\makeatletter
\@ifpackageloaded{tcolorbox}{}{\usepackage[skins,breakable]{tcolorbox}}
\@ifpackageloaded{fontawesome5}{}{\usepackage{fontawesome5}}
\definecolor{quarto-callout-color}{HTML}{909090}
\definecolor{quarto-callout-note-color}{HTML}{0758E5}
\definecolor{quarto-callout-important-color}{HTML}{CC1914}
\definecolor{quarto-callout-warning-color}{HTML}{EB9113}
\definecolor{quarto-callout-tip-color}{HTML}{00A047}
\definecolor{quarto-callout-caution-color}{HTML}{FC5300}
\definecolor{quarto-callout-color-frame}{HTML}{acacac}
\definecolor{quarto-callout-note-color-frame}{HTML}{4582ec}
\definecolor{quarto-callout-important-color-frame}{HTML}{d9534f}
\definecolor{quarto-callout-warning-color-frame}{HTML}{f0ad4e}
\definecolor{quarto-callout-tip-color-frame}{HTML}{02b875}
\definecolor{quarto-callout-caution-color-frame}{HTML}{fd7e14}
\makeatother
\makeatletter
\@ifpackageloaded{caption}{}{\usepackage{caption}}
\AtBeginDocument{%
\ifdefined\contentsname
  \renewcommand*\contentsname{Inhaltsverzeichnis}
\else
  \newcommand\contentsname{Inhaltsverzeichnis}
\fi
\ifdefined\listfigurename
  \renewcommand*\listfigurename{Abbildungsverzeichnis}
\else
  \newcommand\listfigurename{Abbildungsverzeichnis}
\fi
\ifdefined\listtablename
  \renewcommand*\listtablename{Tabellenverzeichnis}
\else
  \newcommand\listtablename{Tabellenverzeichnis}
\fi
\ifdefined\figurename
  \renewcommand*\figurename{Abbildung}
\else
  \newcommand\figurename{Abbildung}
\fi
\ifdefined\tablename
  \renewcommand*\tablename{Tabelle}
\else
  \newcommand\tablename{Tabelle}
\fi
}
\@ifpackageloaded{float}{}{\usepackage{float}}
\floatstyle{ruled}
\@ifundefined{c@chapter}{\newfloat{codelisting}{h}{lop}}{\newfloat{codelisting}{h}{lop}[chapter]}
\floatname{codelisting}{Listing}
\newcommand*\listoflistings{\listof{codelisting}{Listingverzeichnis}}
\makeatother
\makeatletter
\makeatother
\makeatletter
\@ifpackageloaded{caption}{}{\usepackage{caption}}
\@ifpackageloaded{subcaption}{}{\usepackage{subcaption}}
\makeatother

\ifLuaTeX
\usepackage[bidi=basic]{babel}
\else
\usepackage[bidi=default]{babel}
\fi
\babelprovide[main,import]{ngerman}
\ifPDFTeX
\else
\babelfont{rm}[]{Arial}
\fi
% get rid of language-specific shorthands (see #6817):
\let\LanguageShortHands\languageshorthands
\def\languageshorthands#1{}
\ifLuaTeX
  \usepackage[german]{selnolig} % disable illegal ligatures
\fi
\usepackage{bookmark}

\IfFileExists{xurl.sty}{\usepackage{xurl}}{} % add URL line breaks if available
\urlstyle{same} % disable monospaced font for URLs
\hypersetup{
  pdftitle={Lernumgebungen},
  pdfauthor={Richard Conrardy},
  pdflang={de},
  colorlinks=true,
  linkcolor={(3,1.5)},
  filecolor={Maroon},
  citecolor={Blue},
  urlcolor={Blue},
  pdfcreator={LaTeX via pandoc}}


\title{Lernumgebungen}
\author{Richard Conrardy}
\date{}

\begin{document}
\maketitle


\section{Unterrichtsstrukturierung}\label{unterrichtsstrukturierung}

\subsection{Phasenmodelle}\label{phasenmodelle}

Mikroplanung umfasst gemäss der an der PHBern gebräuchlichen Definition
zwei bis vier Lektionen. Die Gliederung des Unterrichts wird jedoch
typischerweise nach anderen ``unterrichtsbestimmende Bauelemente''
(Wollring, 2008, S. 9) aufgeteilt. Diese Aufteilung kann sich auf
Phasenmodelle stützen:

\begin{itemize}
\tightlist
\item
  PADUA: (vgl. Aebli, 1983)
\item
  AVIVA: (Städeli \& Maurer, 2020)
\item
  KAFKA und SAMBA: (Reusser 1999 nach Leuchter, 2009))
\end{itemize}

\begin{tcolorbox}[enhanced jigsaw, title=\textcolor{quarto-callout-tip-color}{\faLightbulb}\hspace{0.5em}{Tipp}, colbacktitle=quarto-callout-tip-color!10!white, breakable, leftrule=.75mm, arc=.35mm, toptitle=1mm, opacityback=0, coltitle=black, left=2mm, opacitybacktitle=0.6, colback=white, bottomtitle=1mm, colframe=quarto-callout-tip-color-frame, titlerule=0mm, rightrule=.15mm, bottomrule=.15mm, toprule=.15mm]

Lernphasen werden in der Lerngelegenheit ``Planungsarbeit im Fokus''
vertieft.

\end{tcolorbox}

\subsection{Lernumgebungen}\label{lernumgebungen}

Eine andere Einteilung ist diejenige in \emph{Lernumgebungen} oder
``Arbeitsumgebung'' (Wollring, 2008, S. 9). Gewisse Autoren
(/Dozierende) setzen diese gleich mit dem Begriff des
\emph{Lernarrangement}, andere definieren sie leicht unterschiedlich.

Van Dormolen (1978) beschreibt eine frühe Form von Lernumgebungen,
nämlich den Unterrichtsarbeitsplan:

\begin{quote}
Ein Unterrichtsarbeitsplan beschreibt die Aktivitäten, von denen Lehrer
und Schüler meinen, dass sie zum Lernziel führen. Er beschreibt weniger
die Stunden selbst als vielmehr die Art der Aktivitäten und wie man die
Ergebnisse testen kann. Ferner werden sowohl der Lehrstoff als auch die
Hilfsmittel, die benutzt werden sollen, ausgewählt.

(Van Dormolen, 1978, S. 1)
\end{quote}

Der Fokus von Unterrichtsarbeitsplänen liegt somit auf Aktivitäten und
die dafür benötigten Hilfsmittel. Lernumgebungen führen diese Begriffe
etwas weiter aus. Eine frühe Definition von \emph{Lernumgebung} findet
man in Friedrich \& Mandl (1997):

\begin{quote}
Unter Lernumgebung wird hier das Arrangement der äußeren Lernbedingungen
(Personen und Institutionen, Geräte und Objekte, Symbole und Medien,
Informationsmittel und Werkzeuge) und Instruktionsmaßnahmen
(Lernaufgaben, Sequenz der Lernschritte, Methoden u. a.) verstanden, die
Lernen ermöglichen und erleichtern.

(Friedrich \& Mandl, 1997, S. 258)
\end{quote}

Eine neuere Definition beschränkt sich auf die wesentlichen vier
Anforderungen.

\begin{quote}
Eine durch den Unterricht hergestellte Lernumgebung besteht aus einem
Arrangement von Unterrichtsmethoden, Unterrichtstechniken,
Lernmaterialien, Medien.

(Reinmann \& Mandl, 2010, S. 615)
\end{quote}

Beide Definitionen lenken die Aufmerksamkeit darauf, dass Lernumgebungen
nicht \emph{one size fits all} erstellt werden, sondern eigentlich die
Realisierung eines spezifischen Unterrichts entsprechen.

``Hier stellt gerade die Grundschule umfassende Anforderungen an das
Implementieren, die in den konfektioniert vorbereiteten
Unterrichtsmaterialien nicht berücksichtigt sein können, sondern von dem
Lehrenden spezifische Entscheidungen und Strukturierungen verlangen, die
der konkreten Situation Rechnung tragen.'' (Wollring, 2008, S. 4)

Eine für den Mathematikunterricht nützliche Sichtweise auf
Lernumgebungen lautet, dass es sich um eine ``flexible grosse Aufgabe''
handelt. Diese wiederum besteht aus einem ``Netzwerk kleinerer
Aufgaben'' (Wollring, 2008, S. 5).

\begin{tcolorbox}[enhanced jigsaw, title=\textcolor{quarto-callout-tip-color}{\faLightbulb}\hspace{0.5em}{Tipp}, colbacktitle=quarto-callout-tip-color!10!white, breakable, leftrule=.75mm, arc=.35mm, toptitle=1mm, opacityback=0, coltitle=black, left=2mm, opacitybacktitle=0.6, colback=white, bottomtitle=1mm, colframe=quarto-callout-tip-color-frame, titlerule=0mm, rightrule=.15mm, bottomrule=.15mm, toprule=.15mm]

An einer Pädagogischen Hochschule (und auch in einem Verlag) können
keine \emph{one size fits all} Materialien erstellt werden. Freie
Lizenzen ermöglichen es jedoch Lehrpersonen Unterrichtsmaterialien für
ihren Unterricht anzupassen.

\end{tcolorbox}

\subsection{Kernfragen}\label{kernfragen}

Damit man bei der Erstellung einer Lernumgebung nicht überwältigt wird,
hilft es sich möglichst früh auf gewisse Eckpfeiler zu einigen.

Das Modell eines Unterrichtsarbeitsplanes, gemäss Van Dormolen (1978),
schlägt folgende vier Kernfragen vor:

\begin{quote}
\begin{itemize}
\tightlist
\item
  Lernziele, ``Was würde ich gerne meine Schüler erreichen lassen?''
\item
  Anfangszustand, ``Wo können meine Schüler anfangen?''
\item
  Unterrichtssituation, ``Wie kann ich Lerninhalte vermitteln und wie
  können meine Schüler sie lernen?''
\item
  Auswertung, ``Wie ermittelt man den Lehrerfolg?''
\end{itemize}

(Van Dormolen, 1978, S. 2--3)
\end{quote}

Diese vier Leitfragen sind selbstverständlich voneinander abhängig und
sie ähneln somit dem \emph{constructive alignment} von Biggs (1996) .

\begin{quote}
\begin{itemize}
\tightlist
\item
  teachers need to be clear about what they want their students to
  learn, and how they would manifest that learning in terms of
  ``performances of understanding''. For .example, memorising and
  paraphrasing are not performances of understanding, recognising an
  application in a novel context is.
\item
  the performance objectives thus emerging need to be arranged in a
  hierarchy from most acceptable to barely satisfactory, which hierarchy
  becomes the grading system.
\item
  students need to be placed in situations that are judged likely to
  elicit the required learnings.
\item
  students are then required to provide evidence, either by self-set or
  teacher-set tasks, as appropriate, that their learning can match the
  stated objectives. Their grade becomes the highest level they can
  match convincingly.
\end{itemize}

(Biggs, 1996, S. 360--361)
\end{quote}

\begin{tcolorbox}[enhanced jigsaw, title=\textcolor{quarto-callout-tip-color}{\faLightbulb}\hspace{0.5em}{Tipp}, colbacktitle=quarto-callout-tip-color!10!white, breakable, leftrule=.75mm, arc=.35mm, toptitle=1mm, opacityback=0, coltitle=black, left=2mm, opacitybacktitle=0.6, colback=white, bottomtitle=1mm, colframe=quarto-callout-tip-color-frame, titlerule=0mm, rightrule=.15mm, bottomrule=.15mm, toprule=.15mm]

Um zu wissen worauf man hinaus möchte, braucht es Erfahrung und Wissen,
welches Sie womöglich noch nicht haben. Holen Sie sich frühzeitig
Feedback!

\end{tcolorbox}

\subsection{Leitideen}\label{leitideen}

Sind die Kernfragen geklärt, kann die Lernumgebung detaillierter geplant
werden. Wollring (2008) schlägt vor die Lernumgebung, als Bündelung
kleiner Aufgaben, an Leitideen auszurichten:

\begin{quote}
\begin{itemize}
\tightlist
\item
  Gegenstand und Sinn, Fach-Sinn und Werk-Sinn
\item
  Artikulation, Kommunikation, Soziale Organisation
\item
  Differenzieren
\item
  Logistik
\item
  Evaluation
\item
  Vernetzung mit anderen Lernumgebungen
\end{itemize}

(Wollring, 2008, S. 5)
\end{quote}

\begin{tcolorbox}[enhanced jigsaw, title=\textcolor{quarto-callout-tip-color}{\faLightbulb}\hspace{0.5em}{Tipp}, colbacktitle=quarto-callout-tip-color!10!white, breakable, leftrule=.75mm, arc=.35mm, toptitle=1mm, opacityback=0, coltitle=black, left=2mm, opacitybacktitle=0.6, colback=white, bottomtitle=1mm, colframe=quarto-callout-tip-color-frame, titlerule=0mm, rightrule=.15mm, bottomrule=.15mm, toprule=.15mm]

Lesen Sie die Ausführungen zu den Leitideen im Originalartikel
aufmerksam durch.

\end{tcolorbox}

\subsection{Unterrichtszentrum}\label{unterrichtszentrum}

Falls Sie den Behaviourismus, Kognitivistismus und Konstruktivismus noch
nicht kennen, lesen Sie Reinmann \& Mandl (2010), ab S.617. Die beiden
ersten Positionen werden als ``technologische Position''
zusammengefasst. Diese Perspektiven können ebenfalls durch ihren Fokus
unterschieden werden: Learner centered, Knowledge centered, Assessment
centered, Community centered (Reinmann \& Mandl, 2010, S. 618)

\phantomsection\label{refs}
\begin{CSLReferences}{1}{0}
\bibitem[\citeproctext]{ref-aebli_zwolf_1983}
Aebli, H. (1983). \emph{Zwölf Grundformen des Lehrens: eine allgemeine
Didaktik auf psychologischer Grundlage: Medien und Inhalte didaktischer
Kommunikation, der Lernzyklus} (15. Auflage). Klett-Cotta.

\bibitem[\citeproctext]{ref-biggs_enhancing_1996}
Biggs, J. (1996). Enhancing teaching through constructive alignment.
\emph{Higher Education}, \emph{32}(3), 347--364.
\url{https://doi.org/10.1007/BF00138871}

\bibitem[\citeproctext]{ref-friedrich_analyse_1997}
Friedrich, H., \& Mandl, H. (1997). Analyse und Förderung
selbstgesteuerten Lernens. In \emph{Lern- Und Denkstrategien} (S.
237--296).

\bibitem[\citeproctext]{ref-leuchter_rolle_2009}
Leuchter, M. (2009). \emph{Die Rolle der Lehrperson bei der
Aufgabenbearbeitung: unterrichtsbezogene Kognitionen von Lehrpersonen}.
Waxmann.

\bibitem[\citeproctext]{ref-krapp_unterrichten_2010}
Reinmann, G., \& Mandl, H. (2010). Unterrichten und Lernumgebungen
gestalten. In A. Krapp \& B. Weidenmann (Hrsg.), \emph{Pädagogische
Psychologie: ein Lehrbuch} (5., vollst. überarb. Aufl., {[}Nachdr.{]}).
Beltz {PVU}.

\bibitem[\citeproctext]{ref-stadeli_aviva_2020}
Städeli, C., \& Maurer, M. (2020). \emph{The {AVIVA} model A
competence-oriented approach to teaching and learning}.
\url{https://www.hep-verlag.ch/the-aviva-model}

\bibitem[\citeproctext]{ref-van_dormolen_didaktik_1978}
Van Dormolen, J. (1978). \emph{Didaktik der Mathematik}. Vieweg+Teubner
Verlag. \url{https://doi.org/10.1007/978-3-322-84149-0}

\bibitem[\citeproctext]{ref-wollring_kennzeichnung_2008}
Wollring, B. (2008). Kennzeichnung von Lernumgebungen für den
Mathematikunterricht in der Grundschule. In \emph{Lernumgebungen auf dem
Prüfstand. Bericht 2 der Kasseler Forschergruppe Empirische
Bildungsforschung Lehren -- Lernen -- Literacy} (S. 9--26). Kasseler
Forschergruppe.
\url{https://www.schulentwicklung.nrw.de/sinus/upload/tagung20080424/2008_Wollring_Lernumgebungen.pdf}

\end{CSLReferences}




\end{document}
